\subsection{Source-proof audit of the direct characteristic-$p$ route}

This audit was carried out on 18 November 2026 by GPT-6 Astra Pro,
Pro mode; a separate reasoning-effort setting was not exposed.
It distinguishes a proof newly supplied here, a source statement
actually opened, and a foundational source theorem used without
claiming an independent proof of its entire theory. The direct
route has the following internal dependencies:
\[
 \mathrm{A.45},\mathrm{A.46},\mathrm{A.47}\Longrightarrow\mathrm{A.48},
 \qquad
 \mathrm{A.45},\mathrm{A.46},\mathrm{A.48}\Longrightarrow\mathrm{A.13}.
\]
The finite-core argument is written out independently inside A.13
and A.48. The former uses A.48 only for its boundary and arithmetic
clause, preserving the inherited statement without a cycle. Thus this
route does not import all of A.3 merely to use that elementary
argument. The replacement A.15 uses A.14 and the actual
compactification and trace-formula arguments written in its proof;
A.49 then uses A.14--A.15 without A.16.

\paragraph{The precise excluded chain.}
The original characteristic-zero route uses SGA~1 III 7.3--7.4
and XIII 2.9--2.11. The alternative citation to Stix's
Proposition 2.7 is valid for relative pro-$r$ completion and
does not itself construct a characteristic-zero lift. However,
Stix's Section 2.3 invokes SGA~1 XIII 4.3--4.4, and the open-pair
example in XIII 4.4 explicitly invokes XIII 2.9. The literal
G1(ii) requirement therefore cannot be met by relabelling that
alternative as a proof free of the excluded source input.
Neither Stix's theorem nor XIII 4.3--4.4 is used in A.45--A.48.

\paragraph{A.45: formal cohomology and the algebraic curve inputs.}
Theorem 12.7 of \citet{CF_Milne2013} is the Leray spectral sequence
for a morphism of schemes and an abelian \emph{\'etale} sheaf. Its
printed proof, pages 82--83, is the derived-functor spectral
sequence for direct image followed by global sections; direct
image preserves injectives. Theorem 14.9, pages 96--97, is the
Cartan--Leray sequence for a finite \emph{\'etale} Galois cover.
Its proof is the same formal construction, using descent and
group invariants. Remark 14.10 passes to continuous cochains on
an inverse tower. These statements contain no specialization or
characteristic-zero comparison. We use them to obtain the two
low-degree Leray sequences and the injection of continuous group
$H^2$ into \emph{\'etale} $H^2$; the latter is proved in A.45 by
the vanishing of $H^1$ on the universal finite-cover tower, not
by presuming the total space to be a $K(\pi,1)$.

Milne's Propositions 14.1--14.2, pages 89--91, provide Picard
divisibility, prime-to-characteristic torsion, and the cohomology
of a proper smooth curve. The proof of Proposition 14.1 first
illustrates the assertion over $\mathbf C$, then gives a separate
algebraic route through the Jacobian. We use the latter route.
To verify that this does not hide a lifting argument, we also
opened \citet[Theorem 8.2 and Remark 8.4]{CF_MilneAV1986} in the
author's corrected 2 January 2022 text, printed pages 13--14.
Theorem 8.2 proves algebraically that $[n]$ has degree $n^{2g}$
and is \emph{\'etale} when $n$ is invertible. Its proof uses a
symmetric ample line bundle, the algebraic theorem of the cube,
intersection numbers, and the differential $d[n]=n$ at the
origin. Remark 8.4 then determines the free rank-$2g$ torsion
module. No comparison with a complex torus is used in that
argument. The author's preface says that the numbering agrees
with the 1986 article; we opened the corresponding original
article PDF as well, but rely for the readable proof on the
corrected author text.

For completeness, the curve-cohomology calculation in Milne
Proposition 14.2 follows from Kummer and Theorem 13.7. We opened
the latter proof, pages 86--88: it uses the Weil-divisor sequence,
Hilbert 90, Tsen's theorem for a function field of one variable
over an algebraically closed field, and the strictly henselian
discrete-valuation version of the algebraic cohomological
dimension calculation. This is an algebraic input valid in
positive characteristic. In A.45, the degree description makes
the pullback on $H^2(S,\mathbf F_r)$ zero under a degree-$r$
\emph{\'etale} cover. Such a cover exists from a nonzero
$r$-torsion Picard class, and the argument remains valid after
each cover because the base genus remains positive.

\paragraph{A.45: purity and proper base change, with open base change derived.}
Milne Theorem 16.7, pages 111--112, is purity for a smooth pair
over the algebraically closed ground field with coefficients of
invertible order. The opened proof reduces \emph{\'etale} locally
to a standard coordinate pair by Lemma 16.8 and starts from the
one-dimensional local calculation of Proposition 14.3. The latter
is the Kummer calculation at a henselian discrete valuation ring.
We use only this smooth-pair theorem; no general absolute-purity
conjecture appearing elsewhere in the old notes is assumed.

Milne Theorem 17.10, pages 116--117, is proper base change for
torsion sheaves and arbitrary base change. It is applied only
to $C_0\times S\to S$ and the proper finite union of labelled
sections. Its proof reduces to the stalk theorem 17.7. The
following source discussion describes the algebraic proof through
the relative Picard functor for curves and induction on relative
dimension, and refers to the detailed proof in Artin's treatment.
We treat proper base change as an established foundational theorem;
we do not claim to have independently reproved that theorem.
The topological discussion at the start of Milne's Section 17 is
an analogy, not an input to the algebraic statement.

The conclusion for the nonproper map $f:U\to S$ is proved in
A.45 from the localization triangle, the proper product and the
proper labelled boundary. In particular the proof does not invoke
smooth base change for $f$, an open homotopy sequence, or the
cohomological properness result SGA~1 XIII 2.9. Fibre compatibility
of localization is checked at $t=0$ by the local Kummer residue
of $t$, which is one and remains one under transverse restriction.
The two constant pieces of $R^1f_*\mathbf F_r$ therefore give an
elementary $r$-group of shears. The sum of the labelled point
classes is a map with every coefficient one; its surjectivity
uses $m\ge1$, not invertibility of $m$.

The remaining nonabelian step is supplied in full in A.45.
For a central subgroup of order $r$, a group-extension factor set
produces the obstruction in continuous $H^2$. The Leray filtration
and finite affine $r$-group representations kill that obstruction
after an $r$-cover of the base. The difference between two lifts
is a central $H^1$ class, which extends after a further $r$-cover.
This induction does not call on any nonabelian base-change theorem.

\paragraph{A.46: checking the proof of the tame-chart source.}
SGA~1 XIII Proposition 5.2 and Corollary 5.3 were opened with
their proof, original margin pages 428--431. The source proof
of 5.2 uses regularity of the Kummer chart from 5.1, uniqueness
of normalization, purity away from codimension one, and the
discrete-valuation Abhyankar lemma SGA~1 X 3.6. We opened the
latter proof, original margin pages 279--280: the composite
inertia embeds in a product of cyclic tame inertia groups; if
one order is divisible by the other, its projection to the
larger cyclic group is an isomorphism. This is a finite cyclic
group argument. The specialization discussion following that
lemma is not used.

The purity cited by XIII 5.2 is SGA~2 XIV 1.11. We downloaded
and opened its actual statement and proof, together with SGA~2
X 3.4, to which it reduces. In the 2005 SMF edition these are
printed pages 169 and 94--97, respectively. The regular case
extends covers across a subset of codimension at least two.
The proof in dimension two extends the finite algebra, uses
depth and the Auslander--Buchsbaum formula to make it locally
free, and then the discriminant ideal excludes ramification
supported only at the closed point. Higher dimension follows
by completion and induction on a regular parameter, with
effective descent. These are algebraic local arguments, with
no characteristic-zero lift and no SGA~1 XIII 2.9--2.11 input.
In particular the relevant purity theorem is the homotopical
theorem XIV 1.11, not the separate cohomological theorem XIV 1.10
on the same page.

The other source input in A.46 is SGA~1 X Proposition 1.2,
original margin page 262, for the finite \emph{\'etale} Stein
factor of a proper separable morphism. Its proof reduces to
a complete local base and the coherent $H^0$ base-change and
flatness argument. Here the tame charts make the normalization
smooth proper over the curve, so its fibres are geometrically
reduced and every hypothesis is met. The source proof cites
EGA~III 7.8.10(i), not an open-family specialization theorem.
The commuting transitive-action argument proving normality of
the fibre image and the quotient from the base is supplied in
A.46 itself.

\paragraph{A.47 and A.48: inertia and boundary descent.}
A.47 uses only the algebraic Picard/Kummer inputs just checked.
It constructs branched cyclic covers of order $r^n$ after a proper
cyclic $r$-cover, and takes their conjugate compositum. Thus the
unbounded order of actual cusp inertia is established without a
surface presentation or lifting to characteristic zero.

For A.48, the language of root stacks is a description of
\emph{\'etale} descent from the charts in SGA~1 XIII 5.2. The
source itself contains no root-stack assertion. Unique extension
extends the transition isomorphisms and their equivariance;
ordinary finite \emph{\'etale} descent glues them. Descent from
the boundary gerbe for covers with trivial stabilizer action gives
the right exact boundary sequence. Its finite stabilizer is not
assumed injective. The inverse limit, right exactness of the
geometrically pro-$\Sigma$ quotient, and the actual transverse
fibre map are written out in A.48. Full infinite inertia follows
from A.47 and the geometric fibre injection proved at the start of
A.48 directly from A.45--A.46. The final two conjugations are
group-theoretic and preserve the full arithmetic quotient.

\paragraph{What the dependency audit establishes.}
The direct proofs have no edge to curve lifting, to the
characteristic-zero configuration theorem, to Stix's exactness
theorem, or to SGA~1 XIII 2.9--2.11. Their foundational inputs
remain algebraic purity, proper base change, Picard theory,
Kummer theory, descent, and formal spectral sequences. This is
a proof-dependency statement, not a claim that every foundational
theorem has been independently reproved here. The unresolved
published-number match for Stix's proposition is therefore a
bibliographic limitation of the alternative route and does not
affect the direct proof or its stated prime-set range.
